% SPDX-License-Identifier: GPL-3.0-or-later
% Spanish (español) -- what the reading editions' preamble leaves to the
% language.  Input by lib/frank-preamble.tex as
% \input{\FrankLib/lang/\BookLang.tex}; docs/languages.md section 6 lists the
% hooks every file here must define and section 12 the ones it may.
% A Latin-script language, the same shape as lib/lang/it.tex: no font of its
% own, no marks to strip, two passes.

% ---- the babel language and the switches ----------------------------------
% \FrankLangName is what \foreignlanguage and \begin{otherlanguage} are given
% everywhere in the core preamble: babel must know this name, or every passage
% comes out in the roman with a "Unknown language" error.
\newcommand{\FrankLangName}{spanish}
\FrankRTLfalse           % chunk left, gloss right; \pw needs no bidi isolate
\FrankHasBarefalse       % nothing comes off, so pass 3 would repeat pass 1
\FrankHasAltfalse        % whether \FrankAltFont exists (the title, the chapter number)
\FrankHasReadingfalse    % whether \chr's third argument is set as a reading
\FrankHasVertfalse       % whether pass 4 is \FrankVertPass rather than the font pass

% ---- fonts ------------------------------------------------------------------
% No face is NAMED here: \FrankPickFont walks the registry's tex.main and then
% tex.main_fallbacks and takes the first one installed, so the .tex and the
% .json cannot drift apart (docs/languages.md section 12).  Spanish names
% none: tex.main is null in the registry, so the text is set in the roman the
% core preamble already loaded -- the same face as the glosses, which is right
% for a Latin script and wrong for any other.  import=es brings babel's
% Spanish hyphenation (the registry's tex.hyphen) and its quoting rules.
\babelprovide[import=es]{spanish}

% ---- the vocabulary line -----------------------------------------------------
% \vb{infinitive}{sound}{first person present}{sound}{third person preterite}{sound}{meaning}
% Those two are the forms a reader cannot guess: the yo present carries the odd
% consonant (tengo, salgo, conozco) and the stem change (quiero, pido), and is
% what the present subjunctive is built on; the él preterite carries a stem of
% its own (tuvo, dijo, fue) and the -ir stem change the yo form hides.
%
% THE THIRD SLOT WAS THE FIRST PERSON PRETERITE, and its label "past".  Both
% were wrong for Spanish.  pedí, sentí, dormí, leí show nothing of what pidió,
% sintió, durmió, leyó show -- the e to i, the o to u, the y -- which is why
% the conventions had to write (e→i) beside pedí; the third person is the
% form a story is told in, and its stem comes back in the imperfect
% subjunctive and the gerund (pidiera, pidiendo).  All the yo preterite adds
% is spelling (llegué, busqué).  And "past" named one of two simple pasts as
% THE past, in a language whose conventions spend a paragraph on the preterite
% and the imperfect not being one tense: pret. is what es.md already wrote
% everywhere else, and what German's pair already says.
%
% An impersonal verb has no yo to give, and takes the third person there too
% (llueve, hay).  What the three forms cannot say goes in ONE parenthesis
% after the meaning, items parted by "; ", and only where it is irregular:
%   p.p. \pw{hecho}     the past participle (dicho, visto, puesto, escrito ...)
%   fut. \pw{tendré}    the future, whose stem also makes the conditional
%   \vb{hacer}{}{hago}{}{hizo}{}{to do, to make (p.p. \pw{hecho}; fut. \pw{haré})}
% Where there are two and one is irregular, both, parted by a slash:
% p.p. \pw{imprimido}/\pw{impreso}.  The \vb the sources sidebar drafts
% (lib/verbs/es.py) writes exactly this, from the dictionary's own rows.
% The class label (-ar, -er, -ir) and the stem-change label (e→i) are gone:
% the infinitive and the two forms show both.
% The pair below comes from the registry's "vb_labels" ("pres.", "pret."), which
% is where the reader takes it from as well, beside "vb_forms", the three forms
% in words: change one and change the other, or the PDF and the reader of one
% book label the same form differently.  docs/lang/es.md names the same three
% forms in the same order.  lib/newlang.py --check compares them.
\newcommand{\FrankVbPres}{pres.}
\newcommand{\FrankVbPast}{pret.}

% ---- the two Lua filters ---------------------------------------------------
% Both are the identity: Spanish carries no marks the bare pass would drop
% (the registry's `strip` is null -- the written accent is spelling and stays
% in every pass) and its labels are already in Latin digits.  They still have
% to exist, because the core calls them for every language.  No #, % or
% literal backslash in a \directlua body: TeX reads it first (NOTES section 9
% of the Persian book).
\directlua{
  function frank_strip(s) return s end
  function frank_latin(s) return s end
}

% ---- the "How to read this" page -----------------------------------------------
% Printed by lib/frank-frontmatter.tex, before the first chapter: the only
% place the reader is told what the passes are for.  Spanish words are set with
% \textit and not \pw, as Italian's page does: the language is in the same
% roman as the prose round it, so there is no face to switch to and \pw would
% only make them smaller.
\newcommand{\FrankHowTo}{%
Every passage comes twice.

\smallskip
\textbf{First} the whole sentence, and nothing else. Try to read it. Get as
far as you can and do not worry about what you cannot get — the point is the
attempt, made before any help arrives.

\smallskip
\textbf{Second} the same sentence in small chunks, Spanish on the left, and
on the right three lines: how it sounds, then how to look each word up, then
what it means. Now you find out how much you had right. Spanish is written as
it is said, so there is no third pass: the sentence you attempted is already
the sentence as it is printed.

\smallskip
The pronunciation line is given only where the spelling misleads, which is
seldom, and never marks stress: a word stressed against the ordinary rule
carries a written accent already (\textit{corazón}). What it shows is the few
letters that do not say themselves — \textit{x} for the \textit{ch} of
\textit{loch} (\textit{j}, and \textit{g} before \textit{e} or \textit{i}); a
silent \textit{h}; \textit{ʝ} for \textit{ll} and \textit{y}; \textit{ɲ} for
\textit{ñ}; \textit{rr} for the trill, which \textit{r} also is at the start
of a word. \textit{c} before \textit{e} or \textit{i}, and \textit{z}, are
\textit{s} in the Americas and \textit{θ} in most of Spain, whichever the
book's own recording says; \textit{b} and \textit{v} are one sound.

\smallskip
Verbs are given as infinitive, the present as \textit{yo} says it, the
preterite as \textit{él} or \textit{ella} says it, meaning — the two middle
forms are the ones you cannot guess (\textit{tener}, \textit{tengo},
\textit{tuvo}; \textit{pedir}, \textit{pido}, \textit{pidió}), and
everything else is built from them. An irregular past participle and an
irregular future stand in brackets after the meaning (\textit{p.p. hecho},
\textit{fut. tendré}): the future's stem makes the conditional too, and
\textit{tendría} is on every page of a story.

\smallskip
Spanish has two verbs for \textit{to be} and the line always says which:
\textit{ser} for what a thing is, \textit{estar} for where and how it is at
this moment. A subjunctive is named where it appears, since your own language
may have no form for that mood.
}
